\begin{tikzpicture}[font=\small,line cap=round,line join=round]\begin{scope}[x=0.75cm,y=0.03333333333333333cm]\draw[otRule!65,line width=.2pt](0,0)--(0,120);\draw[otRule!65,line width=.2pt](2,0)--(2,120);\node[below,font=\tiny,fill=white,inner sep=.5pt]at(2,0){2};\draw[otRule!65,line width=.2pt](4,0)--(4,120);\node[below,font=\tiny,fill=white,inner sep=.5pt]at(4,0){4};\draw[otRule!65,line width=.2pt](6,0)--(6,120);\node[below,font=\tiny,fill=white,inner sep=.5pt]at(6,0){6};\draw[otRule!65,line width=.2pt](8,0)--(8,120);\node[below,font=\tiny,fill=white,inner sep=.5pt]at(8,0){8};\draw[otRule!65,line width=.2pt](10,0)--(10,120);\node[below,font=\tiny,fill=white,inner sep=.5pt]at(10,0){10};\draw[otRule!65,line width=.2pt](12,0)--(12,120);\node[below,font=\tiny,fill=white,inner sep=.5pt]at(12,0){12};\draw[otRule!65,line width=.2pt](0,0)--(12,0);\draw[otRule!65,line width=.2pt](0,20)--(12,20);\node[left,font=\tiny,fill=white,inner sep=.5pt]at(0,20){20};\draw[otRule!65,line width=.2pt](0,40)--(12,40);\node[left,font=\tiny,fill=white,inner sep=.5pt]at(0,40){40};\draw[otRule!65,line width=.2pt](0,60)--(12,60);\node[left,font=\tiny,fill=white,inner sep=.5pt]at(0,60){60};\draw[otRule!65,line width=.2pt](0,80)--(12,80);\node[left,font=\tiny,fill=white,inner sep=.5pt]at(0,80){80};\draw[otRule!65,line width=.2pt](0,100)--(12,100);\node[left,font=\tiny,fill=white,inner sep=.5pt]at(0,100){100};\draw[otRule!65,line width=.2pt](0,120)--(12,120);\node[left,font=\tiny,fill=white,inner sep=.5pt]at(0,120){120};\draw[otMuted,->](0,0)--(12.2,0)node[right]{min};\draw[otMuted,->](0,0)--(0,126.0)node[above]{$T$ ($^\circ$C)};\node[below left,font=\tiny]at(0,0){0};\begin{scope}\clip(0,0)rectangle(12,120);\draw[geoTrait](0,20)--(4,100)--(6,100)--(12,40);\geoSolution{}\end{scope}\end{scope}\end{tikzpicture}\par\vspace{1mm}{\fontsize{8}{10}\selectfont\color{otMuted}Les segments du modèle relient les températures indiquées. Les axes sont gradués ; leurs échelles peuvent différer.}